\documentclass[11pt,a4paper]{article}
\usepackage[T1]{fontenc}
\usepackage[{{babel}}]{babel}
\usepackage{lmodern}
\usepackage[margin=2cm]{geometry}
\usepackage{amsmath, amssymb}
\usepackage{tcolorbox}
\usepackage{tikz}

{{fr:% Sujet et corrigé : \corrigetrue montre les corrigés, \corrigefalse les cache
% (l'aperçu de RayTeX bascule cette ligne, et exporte le sujet et le corrigé).
\newif\ifcorrige
\corrigefalse

% Le corrigé : encadré dans le corrigé, absent du sujet.
\tcbuselibrary{breakable}
\AtBeginDocument{%
  \ifcorrige
    \newtcolorbox{corrige}{breakable, colback=green!3, colframe=green!45!black, fonttitle=\bfseries, title=Corrigé}%
  \else
    \NewDocumentEnvironment{corrige}{+b}{}{}%
  \fi}

% Un exercice numéroté, avec son barème : \begin{exercice}[4 points].
\newcounter{exercice}
\NewDocumentEnvironment{exercice}{o}
  {\par\bigskip\refstepcounter{exercice}\noindent\textbf{Exercice~\theexercice}\IfValueT{#1}{\enspace\textit{(#1)}}\par\nobreak\smallskip}
  {\par}

% Des lignes pour répondre, dans le sujet seulement : \lignesreponse{5}.
\newcounter{lignesreponse}
\newcommand\lignesreponse[1]{\ifcorrige\else
  \par\setcounter{lignesreponse}{#1}%
  \loop\ifnum\value{lignesreponse}>0
    \vspace{4mm}\noindent\makebox[\linewidth]{\dotfill}\par
    \addtocounter{lignesreponse}{-1}%
  \repeat
\fi}

% Une place vide pour répondre, dans le sujet seulement : \espacereponse{3cm}.
\newcommand\espacereponse[1]{\ifcorrige\else
  \par\vspace*{#1}
\fi}

% Un quadrillage de 5 mm pour répondre, dans le sujet seulement : \quadrillage{4cm}.
\newcommand\quadrillage[1]{\ifcorrige\else
  \par\medskip\noindent
  \begin{tikzpicture}\draw[step=5mm, gray!40, very thin] (0,0) grid (\linewidth,#1);\end{tikzpicture}\par
\fi}|en:% Subject and answer key: \answerstrue shows the solutions, \answersfalse hides them
% (the preview of RayTeX turns this line, and exports both).
\newif\ifanswers
\answersfalse

% A solution: framed in the answer key, left out of the subject.
\tcbuselibrary{breakable}
\AtBeginDocument{%
  \ifanswers
    \newtcolorbox{solution}{breakable, colback=green!3, colframe=green!45!black, fonttitle=\bfseries, title=Solution}%
  \else
    \NewDocumentEnvironment{solution}{+b}{}{}%
  \fi}

% A numbered exercise, with its points: \begin{exercise}[4 points].
\newcounter{exercise}
\NewDocumentEnvironment{exercise}{o}
  {\par\bigskip\refstepcounter{exercise}\noindent\textbf{Exercise~\theexercise}\IfValueT{#1}{\enspace\textit{(#1)}}\par\nobreak\smallskip}
  {\par}

% Lines to answer on, in the subject only: \answerlines{5}.
\newcounter{answerlines}
\newcommand\answerlines[1]{\ifanswers\else
  \par\setcounter{answerlines}{#1}%
  \loop\ifnum\value{answerlines}>0
    \vspace{4mm}\noindent\makebox[\linewidth]{\dotfill}\par
    \addtocounter{answerlines}{-1}%
  \repeat
\fi}

% Room to answer, in the subject only: \answerspace{3cm}.
\newcommand\answerspace[1]{\ifanswers\else
  \par\vspace*{#1}
\fi}

% A 5 mm grid to answer on, in the subject only: \answergrid{4cm}.
\newcommand\answergrid[1]{\ifanswers\else
  \par\medskip\noindent
  \begin{tikzpicture}\draw[step=5mm, gray!40, very thin] (0,0) grid (\linewidth,#1);\end{tikzpicture}\par
\fi}}}

\begin{document}

\noindent\textbf{{{institution}}} \hfill {{fr:Durée : 1 h|en:Duration: 1 h}}
\begin{center}
  {\Large\bfseries {{title}}}
\end{center}
\noindent {{fr:Nom|en:Name}}~: \rule{5cm}{0.4pt} \hfill {{fr:Classe|en:Class}}~: \rule{3cm}{0.4pt}

\begin{{{fr:exercice|en:exercise}}}[4 points]
  {{fr:Résoudre dans $\mathbb{R}$ l'équation $x^2 - 5x + 6 = 0$.|en:Solve $x^2 - 5x + 6 = 0$ over $\mathbb{R}$.}}
  \{{fr:lignesreponse|en:answerlines}}{4}
  \begin{{{fr:corrige|en:solution}}}
    $\Delta = 25 - 24 = 1$, {{fr:donc|en:so}} $x = \frac{5 - 1}{2} = 2$ {{fr:ou|en:or}} $x = \frac{5 + 1}{2} = 3$.
  \end{{{fr:corrige|en:solution}}}
\end{{{fr:exercice|en:exercise}}}

\begin{{{fr:exercice|en:exercise}}}[3 points]
  {{fr:Placer les points $A(1, 2)$ et $B(3, -1)$ dans un repère, puis tracer la droite $(AB)$.|en:Plot the points $A(1, 2)$ and $B(3, -1)$, then draw the line $(AB)$.}}
  \{{fr:quadrillage|en:answergrid}}{4cm}
  \begin{{{fr:corrige|en:solution}}}
    {{fr:La droite $(AB)$ a pour coefficient directeur|en:The slope of $(AB)$ is}} $\frac{-1 - 2}{3 - 1} = -\frac{3}{2}$.
  \end{{{fr:corrige|en:solution}}}
\end{{{fr:exercice|en:exercise}}}

\begin{{{fr:exercice|en:exercise}}}[3 points]
  {{fr:Énoncer le théorème de Pythagore.|en:State the Pythagorean theorem.}}
  \{{fr:espacereponse|en:answerspace}}{3cm}
  \begin{{{fr:corrige|en:solution}}}
    {{fr:Dans un triangle rectangle, le carré de l'hypoténuse est égal à la somme des carrés des deux autres côtés.|en:In a right triangle, the square of the hypotenuse equals the sum of the squares of the other two sides.}}
  \end{{{fr:corrige|en:solution}}}
\end{{{fr:exercice|en:exercise}}}

\end{document}
